% Detailed original version of the ANN and LSTM theory section.
% This file is intentionally not included by main.tex; it preserves the longer mathematical treatment.

\section{Artificial Neural Networks}
\label{sec:anns}
\ac{ANN}s have become widely known, in particular not later than the public release of ChatGPT in 2022 \cite{openai_introducing_2024}. However, the underlying concepts and model principles have existed for decades. Their limited relevance in earlier years was due to several constraints, including data scarcity, hardware limitations, less effective optimisation techniques, and the lack of sufficiently powerful architectures \cite{wu_ai_2025}. Today, these limiting factors have largely been overcome.

\ac{ANN}s, or \ac{NN}s, are a class of machine learning models inspired by biological brains. A neural network consists of interconnected computational units, called neurons, that are organised into layers. Its basic building block is the perceptron, which computes a weighted sum of its inputs, adds a bias term, and applies a non-linear activation function:

\begin{equation}
  y = \sigma \left( \sum_{j=1}^{d} w_j x_j + b \right) = \sigma(\mathbf{w}^{\top} \mathbf{x} + b)
  \label{eq:perceptron}
\end{equation}

where $\mathbf{x} \in \mathbb{R}^{d}$ is the input vector, $\mathbf{w} \in \mathbb{R}^{d}$ is the weight vector, $b$ is the bias term, and $\sigma(\cdot)$ is a non-linear activation function such as the sigmoid, hyperbolic tangent ($\tanh$), or Rectified Linear Unit (ReLU).

A feedforward neural network, also called a multilayer perceptron, consists of multiple layers of neurons, where the output of one layer serves as the input to the next. For a network with $L$ layers, the forward computation can be written recursively as

\begin{align}
  \mathbf{h}^{(0)} &= \mathbf{x} \\
  \mathbf{h}^{(l)} &= \sigma^{(l)} \left( \mathbf{W}^{(l)} \mathbf{h}^{(l-1)} + \mathbf{b}^{(l)} \right), \quad l = 1, \ldots, L
  \label{eq:feedforward-network}
\end{align}

where $\mathbf{h}^{(l)}$ denotes the activations of layer $l$, $\mathbf{W}^{(l)}$ and $\mathbf{b}^{(l)}$ are the weight matrix and bias vector of layer $l$, and $\sigma^{(l)}$ is the activation function applied in that layer.

\subsection{Training Neural Networks}
\label{subsec:training_nns}
Training a neural network is finding the values of its weights and biases that minimise a loss function. This optimisation problem is commonly solved by gradient descent, which updates the parameter vector $\boldsymbol{\theta}$ iteratively in the direction of steepest descent:

\begin{equation}
  \boldsymbol{\theta}_{t+1} = \boldsymbol{\theta}_{t} - \eta \, \nabla_{\boldsymbol{\theta}} \mathcal{L}(\boldsymbol{\theta}_{t})
  \label{eq:gradient-descent}
\end{equation}

where $\eta > 0$ is the learning rate and $\nabla_{\boldsymbol{\theta}} \mathcal{L}$ is the gradient of the loss function with respect to the model parameters. In neural networks, these gradients are computed efficiently by backpropagation, which applies the chain rule of calculus layer by layer from the output back to the input. For layer $l$, the gradient with respect to the weight matrix can be expressed as

\begin{equation}
  \frac{\partial \mathcal{L}}{\partial \mathbf{W}^{(l)}} = \frac{\partial \mathcal{L}}{\partial \mathbf{h}^{(l)}} \cdot \frac{\partial \mathbf{h}^{(l)}}{\partial \mathbf{W}^{(l)}}
  \label{eq:backprop-weight-gradient}
\end{equation}

and the gradient with respect to the previous layer is propagated recursively according to

\begin{equation}
  \frac{\partial \mathcal{L}}{\partial \mathbf{h}^{(l-1)}} = \left( \mathbf{W}^{(l)} \right)^{\top} \frac{\partial \mathcal{L}}{\partial \mathbf{z}^{(l)}} \odot \sigma'^{(l-1)} \left( \mathbf{z}^{(l-1)} \right),
  \label{eq:backprop-hidden-gradient}
\end{equation}

where $\mathbf{z}^{(l)} = \mathbf{W}^{(l)} \mathbf{h}^{(l-1)} + \mathbf{b}^{(l)}$ denotes the pre-activation of layer $l$, and $\odot$ denotes element-wise multiplication. In practice, computing gradients over the entire training set is computationally expensive. For this reason, training is usually performed with mini-batch-based variants such as \ac{SGD}, which estimate the gradient from a small random subset of the data. Many practical training algorithms build on this principle, including momentum-based methods and adaptive optimisers such as AdaGrad, RMSprop, and Adam.

Neural networks are particularly well suited for learning complex, non-linear input--output relationships from large amounts of data. This makes them highly effective in tasks such as perception, prediction, function approximation, and control, especially when manually designing suitable features is difficult. At the same time, they typically require substantial amounts of training data, careful hyperparameter tuning, and considerable computational resources. In addition, their internal representations are often difficult to interpret, and their performance can degrade when training data are unrepresentative or when the optimisation problem is poorly conditioned. Nevertheless, their flexibility and expressive power make them a central modelling approach in modern machine learning. In particular, recurrent architectures extend these ideas to sequential data, which motivates the discussion of long short-term memory networks in the next subsection.

\subsection{Long Short-Term Memory}
Long short-term memory (LSTM) networks are a special type of \ac{RNN} introduced by Hochreiter and Schmidhuber (1997) to model sequential data while adressing the vanishing-gradient problem of conventional RNNs \cite{hochreiter_long_1997}. In contrast to feedforward networks, which process each input independently, recurrent architectures maintain an internal state that is updated over time. This makes them suitable for tasks in which the current output depends not only on the present input but also on those before. Such temporal dependencies are highly relevant in many engineering systems, including hybrid electric vehicles, where current system behaviour depends on past operating conditions, driver actions, and internal states.

The key idea of an LSTM is the introduction of a memory cell whose content can be preserved, overwritten, or exposed in a controlled manner. This is achieved by a set of gates that regulate the flow of information through the network. For an input vector $\mathbf{x}_t$, previous hidden state $\mathbf{h}_{t-1}$, and previous cell state $\mathbf{c}_{t-1}$ at time step $t$, the LSTM updates are commonly written as

\note{check the formulae again}

\begin{align}
  \mathbf{f}_t &= \sigma\left(\mathbf{W}_f \mathbf{x}_t + \mathbf{U}_f \mathbf{h}_{t-1} + \mathbf{b}_f\right) \\
  \mathbf{i}_t &= \sigma\left(\mathbf{W}_i \mathbf{x}_t + \mathbf{U}_i \mathbf{h}_{t-1} + \mathbf{b}_i\right) \\
  \tilde{\mathbf{c}}_t &= \tanh\left(\mathbf{W}_c \mathbf{x}_t + \mathbf{U}_c \mathbf{h}_{t-1} + \mathbf{b}_c\right) \\
  \mathbf{c}_t &= \mathbf{f}_t \odot \mathbf{c}_{t-1} + \mathbf{i}_t \odot \tilde{\mathbf{c}}_t \\
  \mathbf{o}_t &= \sigma\left(\mathbf{W}_o \mathbf{x}_t + \mathbf{U}_o \mathbf{h}_{t-1} + \mathbf{b}_o\right) \\
  \mathbf{h}_t &= \mathbf{o}_t \odot \tanh\left(\mathbf{c}_t\right)
  \label{eq:lstm}
\end{align}

where $\mathbf{f}_t$ is the forget gate, $\mathbf{i}_t$ the input gate, $\mathbf{o}_t$ the output gate, $\tilde{\mathbf{c}}_t$ the candidate cell state, and $\odot$ denotes element-wise multiplication. The forget gate determines which information from the previous cell state should be retained, the input gate controls how much new information is written into memory, and the output gate decides which part of the updated cell state is exposed as hidden state. Through this gated structure, LSTMs can preserve relevant information over longer time horizons than standard RNNs while keeping the gradients from vanishing.

% \hl{This ability makes LSTMs particularly useful for time-series regression, sequence prediction, and system modelling problems in which delayed effects and temporal correlations are important. In the context of this thesis, these properties are especially relevant because the behaviour of a hybrid powertrain is not determined by the instantaneous input alone. Instead, it depends on temporal patterns such as previous power requests, battery state evolution, and drivetrain dynamics. LSTM-based surrogate models can therefore capture dynamic relationships more effectively than purely static feedforward models. At the same time, they remain substantially cheaper to evaluate than high-fidelity physics-based simulations, which makes them attractive for data-intensive downstream tasks such as reinforcement learning.}
\note{commented out: nice tho, but since i am not using a recurrent agent, its kinda hypocritic; maybe leave out or adjust or explicitly mention; probably leave out since this chapter should only state info in general}
